\chapter*{Notation and Conventions}
\addcontentsline{toc}{chapter}{Notation and Conventions}

The following mathematical conventions and notational standards are maintained
throughout the manuscript.

\section*{Fundamental Units and Geometry}
Natural units are employed with $c = \hbar = 1$.
Angles are measured relative to the circle constant $\tau$, defined as the ratio
of circumference to radius, satisfying $\tau = 2\pi$.
Spacetime is modelled as a four-dimensional affine manifold $\mathcal{M}$
equipped with a non-degenerate symmetric bilinear form $g$ of Lorentzian
signature $(1, 3)$.

The Einstein summation convention applies to all repeated indices.
Greek indices $\mu, \nu$ range over spacetime dimensions $0, 1, 2, 3$.
Latin indices $i, j, k$ denote spatial components $1, 2, 3$.
The Minkowski metric components evaluated in a standard orthonormal basis are
$\eta_{\mu\nu} = \mathrm{diag}(1, -1, -1, -1)$.

\section*{Clifford Algebra}

A general real quadratic space $(V, g)$ is characterised by the signature
triplet $(p, m, z)$, where $p, m, z$ denote the number of orthogonal basis
elements squaring to $+1, -1, 0$, respectively.
The associated Clifford algebra is denoted by $\mathcal{C}\ell_{p,m,z}(V, g)$,
and its even subalgebra is denoted by $\mathcal{C}\ell^+_{p,m,z}$.
When the metric is non-degenerate, $z = 0$, these algebras are written
compactly as $\mathcal{C}\ell_{p,m}$ and $\mathcal{C}\ell^+_{p,m}$.

The spacetime algebra $\mathcal{C}\ell_{1,3}$ is generated by an orthonormal
basis of vectors $\{\gamma_\mu\}$ satisfying the anticommutation relation
$\gamma_\mu \gamma_\nu + \gamma_\nu \gamma_\mu = 2\eta_{\mu\nu}$.
For an arbitrary multivector $A$, the operation $\grade{A}{k}$ denotes
projection onto the grade-$k$ subspace.
The unit pseudoscalar is $I = \gamma_0 \gamma_1 \gamma_2 \gamma_3$, satisfying
$I^2 = -1$.

\section*{Algebraic Operations and Involutions}
The primary binary operations between multivectors $A$ and $B$ of pure grades
$r$ and $s$ are defined as follows:

\begin{table}[H]
  \centering
  \begin{tabular}{ll}
    \toprule
    Operation & Notation and Definition \\
    \midrule
    Geometric product & $A B$ \\
    Exterior product & $A \wedge B = \grade{AB}{r+s}$ \\
    Left contraction & $A \lc B = \grade{AB}{s-r} \quad (s \ge r)$ \\
    Right contraction & $A \rc B = \grade{AB}{r-s} \quad (r \ge s)$ \\
    Anticommutator & $A \circ B = \frac{1}{2}(AB + BA)$ \\
    Commutator & $A \diamond B = \frac{1}{2}(AB - BA)$ \\
    \bottomrule
  \end{tabular}
\end{table}

The algebra admits three intrinsic linear operations:
\begin{itemize}
  \item Grade involution $\hat{A}$: Reverses the sign of odd-graded components,
    $\hat{A} = \sum_k (-1)^k \grade{A}{k}$.
  \item Reversion $\tilde{A}$: Reverses the order of constituent vector products,
    $\tilde{A} = \sum_k (-1)^{k(k-1)/2} \grade{A}{k}$.
  \item Clifford conjugation $\hat{\tilde{A}}$: The composition of grade
    involution and reversion.
\end{itemize}

Relative to a chosen observer with timelike basis vector $\gamma_0$, two
frame-dependent operations are defined:
\begin{itemize}
  \item Spatial grade involution: $\bar{A} = \gamma_0 A \gamma_0$.
  \item Frame reversal (Dirac conjugate): $A^\dagger = \gamma_0 \tilde{A}
    \gamma_0$.
\end{itemize}

\section*{Transformation Groups}
The principal transformation groups governing relativistic spacetime symmetries
are defined as follows:

\begin{itemize}
  \item $\mathrm{O}(1,3)$: The Lorentz group, comprising all linear isometries
    preserving the spacetime metric $g$.
  \item $\mathrm{IO}(1,3)$: The Poincaré group (inhomogeneous Lorentz group),
    defined by the semi-direct product $\mathrm{T}(1,3) \rtimes \mathrm{O}(1,3)$
    of spacetime translations and Lorentz transformations.
  \item $\Gamma(V, g)$: The Clifford (Lipschitz) group, consisting of
    invertible multivectors $R \in \mathcal{C}\ell^*(V, g)$ mapping $V \to V$
    under the sandwich action $v \mapsto R v R^{-1}$.
  \item $\mathrm{Spin}^+(1,3)$: The restricted spin group, the double cover of
    the proper orthochronous Lorentz group $\mathrm{SO}^+(1,3)$, comprising
    even multivectors $R \in \mathcal{C}\ell^+_{1,3}$ satisfying $R\tilde{R} =
    1$.
\end{itemize}

\section*{Figures and Visual Representations}
Unless accompanied by an explicit bibliographic citation, all geometric
diagrams, spacetime projections, and figures throughout this monograph
are original illustrations constructed by the author.
